\chapter{The remaining arithmetic theorem and a verifiable research loop}
\label{book:frontier}
\section{A precise positive spectral target}
For a primitive nonprincipal Dirichlet character $\chi$, choose the completed $L$-function with its parity, conductor and gamma factor fixed. A sufficient spectral program requires holomorphic functions $F_N$ on a connected domain containing the critical strip, compact-uniform convergence $F_N\to\Lambda(s,\chi)$, and a proved zero-location property for every $F_N$. A self-adjoint determinant representation must include its domain, determinant regularization, nowhere-zero normalization and coefficient comparison. Self-adjointness of an unrelated finite matrix does not establish any of these identifications.

If every $F_N$ is zero-free off the critical line, Hurwitz on each connected off-line component shows that the nonzero holomorphic limit is zero-free there. The limit cannot be identically zero on an open component if it is a nonzero holomorphic function on the connected containing domain. This last connected-domain hypothesis matters: componentwise convergence alone does not rule out an identically zero limit on a disconnected component. The conclusion is conditional until the arithmetic approximants and positive representation are constructed. General Artin functions also need their own holomorphy statement; Brauer quotients do not supply it. Hecke, Dedekind and elliptic-curve normalization and critical centers must remain distinct.

The original modular-geodesic argument fails at its prime identification. Integer hyperbolic traces and $e^{\ell/2}+e^{-\ell/2}=\sqrt p+1/\sqrt p$ cannot match all rational primes. An asymptotic similarity between prime and geodesic counts does not give a termwise equality of Euler factors. Nor does finite-field purity for a family of curves prove number-field GRH. The local Jacobi theta construction is a positive arithmetic benchmark with its own coefficient and symmetry obstructions, treated in the earlier chapters.

\section{A surviving short-row reduction}
Write $\mathcal O=\Z[\omega]$. For squarefree primary column ideals $n$ of norm about $D$, good away from a fixed set containing 2 and 3, let $\chi_n$ be the original sextic residue character extended by zero at nonunits. At the unit ideal take $\chi_1(u)=1$, including $u=0$. For row-independent coefficients $|c_n|\le1$, the normalized Gram expansion is
\[
 \mathcal E=D^{-1}\sum_{n,m}c_n\overline{c_m}K_H(n,m),\qquad
 K_H(n,m)=\sum_{0<Nu\le H}\chi_n(u)\overline{\chi_m(u)}.
\]
For $n=ha,m=hb$ with $h,a,b$ squarefree and pairwise coprime,
\[
 \chi_n(u)\overline{\chi_m(u)}
 =\mathbf1_{(u,h)=1}\psi_{a,b}(u),\qquad \psi_{a,b}=\chi_a\overline{\chi_b}.
\]
The oscillatory primitive conductor is $ab$; its zero-preserving periodic modulus is $hab$. Cancelling $h$ without its coprimality mask would change the row family. M\"obius inversion gives
\[
 K_H(ha,hb)=\sum_{d\mid h}\mu(d)\psi_{a,b}(d)
  \sum_{0<Nv\le H/Nd}\psi_{a,b}(v).
\]
A row weight has to be evaluated at $dv$ in this identity.

For a fixed nonnegative radial Schwartz profile $\Phi$, replace the sharp row cutoff by $\Phi(u/\sqrt H)$. For an off-diagonal pair, $\psi_{a,b}$ is primitive nonprincipal. If $Q=N(ab)$, its complete Gauss transform and lattice Poisson summation give, for every $X>0$,
\[
 \left|\sum_{v\in\mathcal O}\psi_{a,b}(v)\Phi(v/\sqrt X)\right|\ll_\Phi\sqrt Q.
\]
Indeed the nonzero dual frequencies have Gauss modulus at most $\sqrt Q$ and the zero frequency vanishes. After scaling, Fourier decay bounds their sum by $O_\Phi(Q/X)$: for $X\le Q$ use planar annuli; for $X>Q$ use a decaying power greater than two. The Poisson prefactor is $X/\sqrt Q$. The ideal generator rotates and scales the Eisenstein lattice, so radiality preserves the profile in the divisor sum. This includes scales $X<1$.

For $Nh\asymp G$, the mask costs $\tau(h)\ll_\varepsilon D^\varepsilon$, each kernel is $O_\varepsilon(D^{1+\varepsilon}/G)$, and there are $O(D^2/G)$ pairs. Thus the off-diagonal normalized contribution satisfies
\[
 |\mathcal E_G^{\rm off}|\ll_{\Phi,\varepsilon}\frac{D^{2+\varepsilon}}{G^2}.
\]
The diagonal is $O_\Phi(H)$ for $H\ge1$. At $H=D^{3/4}$, the dyadic range $G\ge D^{5/8}$ is therefore controlled by $O_{\Phi,\varepsilon}(HD^\varepsilon)$. The small-gcd range, especially $h=1$, remains open. This result is for one fixed smooth radial profile; it does not prove the source's stronger moving-profile or supremum uniformity. The derivation and surviving alternative routes are recorded in \texttt{research/short-row-high-gcd.md}.

% Generated from the shared arithmetic article fragments.
\section{Why an absolute coprime pair-square estimate cannot close the short rows}
\label{sec:pair-barrier}
Keep the original sextic family
\[
 M_u(D)=D^{-1/2}\sum_n\mu(n)\nu(n)W(Nn/D)\chi_n(u),
 \qquad \chi_n(u)=(u/n)_6,
\]
with squarefree primary column ideals outside the fixed excluded set,
their full annular coefficients (fix $W$ supported in $[1,2]$), and all
nonzero Eisenstein element rows
$Nu\le H$. The cubic symbol is $\chi_n^2$, not $\chi_n$.
The following classical frame inequality \cite{WelchGeometry} gives a new
obstruction to one proposed estimate in this arithmetic program.

\begin{proposition}[Coprime pair-square obstruction]\label{prop:pair-barrier}
Let $A$ be an $R$-by-$N$ complex matrix with $|A_{ui}|=1$, and put
$K_{ij}=\sum_u A_{ui}\overline{A_{uj}}$. Then
\[
 \sum_{i\ne j}|K_{ij}|^2\ge NR(N-R).
\]
For the original character kernels, let $B_1(H,D)$ be the sum of
$|K_H(n,m)|^2$ over all distinct coprime squarefree column ideals in
$[D,2D]$. At $H=D^{3/4}$,
\[
 B_1(H,D)\gg \frac{D^{11/4}}{(\log D)^2}.
\]
In particular the proposed sufficient input
$B_1(H,D)\ll H^2D^\varepsilon$ fails for each fixed $\varepsilon<5/4$.
This does not refute the signed M\"obius mean square.
\end{proposition}
\begin{proof}
Set $F_{uv}=\sum_i A_{ui}\overline{A_{vi}}$. Reordering four finite sums
proves $\sum_{i,j}|K_{ij}|^2=\sum_{u,v}|F_{uv}|^2$.
The row diagonal contributes $RN^2$; the column diagonal contributes
$NR^2$. Removing the latter and retaining the former gives the inequality.

Restrict the column pairs in the nonnegative sum $B_1$ to distinct prime
ideals $p,q$ of norm in $[D,2D]$, outside the same fixed set. If $H<D$,
no such prime divides a nonzero row element: divisibility would imply
$Nu\ge Np$. Thus all original zero masks are exactly one on this ball,
$|\chi_p(u)|=1$, and the residual primitive finite character conductor is
$pq$, of norm between $D^2$ and $4D^2$. All these pairs have common divisor1.
The lattice count gives $R=(2\pi/\sqrt3)H+O(\sqrt H+1)$.
The classical unconditional fixed-field prime-ideal theorem gives
$N\sim D/\log D$; this is an external analytic input, not a GRH hypothesis.
At the stated endpoint $R=o(N)$, and the bound is therefore
$(1-o(1))N^2R\gg D^{11/4}/(\log D)^2$.
\end{proof}

To see the exact failed step, write $c_n=\mu(n)\nu(n)W(Nn/D)$ and
$C_D=\sum_n|c_n|^2=O(D)$. Cauchy over coprime pairs bounds the absolute
signed contribution by $D^{-1}C_DB_1^{1/2}$. The hoped-for pair-square
upper bound would have closed this particular route, but is false.
A large absolute pair-square mass is not a lower bound on the quadratic
form with the fixed $c_n$. Nor is the prime part a positive summand of the
full M\"obius energy: prime--composite cross terms can cancel. The original
low-gcd estimate and the positive arithmetic determinant bridge remain open.

\subsection{Weighted profiles, exact arithmetic examples and the surviving task}
The finite matrix proof also retains an arbitrary row profile. Suppose
$|A_{ui}|^2=w_u$ for every column, set $S=\sum_u w_u$ and
$Q=\sum_u w_u^2$. The row diagonals are $Nw_u$ and column diagonals are
$S$, so the same identity proves
\[
 \sum_{i\ne j}|K_{ij}|^2\ge N^2Q-NS^2.
\]
This is a finite statement. Applying it to an infinite Schwartz tail would
require its convergence and every newly encountered prime collision;
these are not discarded. If a weighted pair-square attempt includes the
annular $W$, one can instead choose a closed subannulus on which a nonzero
$|W|$ is bounded below. Its prime count still has order $D/\log D$.
The fixed $\nu$ has unit modulus outside its conductor. These facts do not
replace the original coefficients in a signed estimate.

Here are exact sextic computations with all nonzero elements of the
Eisenstein norm ball, and both prime ideals above each split rational prime
in the column annulus. These kernel diagnostics use $\nu=1$ and exclude
only primes above2,3; they do not compute a signed M\"obius energy.
Values are integer norms in $\Z[\omega]$.
\begin{center}\small
\begin{tabular}{rrrrrr}
$D$&$H$&$R$&$N$&$\sum_{i\ne j}|K_{ij}|^2$&$NR(N-R)$\\\hline
128&4&18&22&6840&1584\\
1024&9&36&138&651672&506736\\
4096&81&300&456&60080328&21340800\\
256&64&240&40&363672&$-1920000$
\end{tabular}
\end{center}
The final row has $H=D^{3/4}$ but $N<R$; its negative lower bound is
correctly vacuous. The asymptotic theorem follows from the proof, not an
extrapolation of these examples. For $p\equiv1\pmod3$ and
$w^2+w+1=0$ in $\mathbf F_p$, evaluate $a+b\omega$ as $a+bw$ and take
its $(p-1)/6$ power, lifting the six roots via $\zeta_6=1+\omega$.
Every zero is tested before a unit value is assigned. For instance
$1+3\omega$ has norm7 and vanishes at $\omega=2$ modulo7, so the
unimodularity hypothesis cannot be used beyond its stated domain.

\texttt{ShortRowPairBarrier.lean} verifies the finite frame identity,
row-diagonal bound, weighted inequality and unimodular specialization,
with actual complex sums and explicit hypotheses. It does not formalize
the ideal-character comparison or prime-ideal theorem. The deterministic
integer checker also verifies weighted profiles and256 scalar pair sums
independently of its matrix computation. No numerical model or learned
coefficient is used.

The next signed route must retain $\mu,\nu,W$ before Cauchy. An exact
starting identity in ideal Dirichlet convolution is
\[
 \mu=2a-a*a*\mathbf1+b*b*\mathbf1,\qquad
 a=\mu\,\mathbf1_{Nn\le U},\quad b=\mu-a.
\]
Indeed $\mu*\mathbf1=\delta$ implies
$(\mu-a)*(\mu-a)*\mathbf1=\mu-2a+a*a*\mathbf1$.
On squarefree $n$, each divisor factorization retains pairwise
coprimality; the original row zeros still multiply along it. Estimating
these type-I/type-II kernels with their true conductors is a new open
arithmetic task. The identity alone supplies no cancellation or GRH proof.

% Generated from the shared signed arithmetic fragments.
\section{A signed M\"obius decomposition and its termwise obstruction}
\label{sec:mobius-split}
Keep the original squarefree good-ideal family, its fixed $\nu$, its
full $W(Nn/D)$ and every zero-extended sextic symbol. In ideal Dirichlet
convolution, write $\zeta_0(n)=1$ for every good ideal and let $\delta$
be the convolution unit. For $1\le U<D$ put
$a_U=\mu\mathbf1_{Nn\le U}$ and $b_U=\mu-a_U$. The classical Vaughan
algebra gives
\begin{equation}\label{eq:mobius-split}
 \mu=2a_U-a_U*a_U*\zeta_0+b_U*b_U*\zeta_0.
\end{equation}
Indeed $\mu*\zeta_0=\delta$, and expansion of the last product proves the
identity without an analytic assumption. This is a classical decomposition,
not a new cancellation theorem; compare \cite[Lemma4.1]{GreenTaoVaughan}.

On squarefree $n$ all factors in $n=d c=a b c$ are pairwise coprime.
Define the actual, nonarbitrary coefficient functions
\[
 \beta_U(d)=\sum_{ab=d\atop Na,Nb\le U}\mu(a)\mu(b),\qquad
 g_U(s)=\sum_{b\mid s\atop Nb>U}\mu(b).
\]
For squarefree $d$, $\beta_U(d)$ is $\mu(d)$ times the number of the
displayed ordered factorizations. For $Ns>U$,
$g_U(s)=-\sum_{b\mid s,\,Nb\le U}\mu(b)$. The two regrouped pieces are
\begin{align*}
 I_u&=D^{-1/2}\sum_{d,c\ {\rm squarefree}\atop(d,c)=1}
  \beta_U(d)\nu(dc)\chi_d(u)\chi_c(u)W(NdNc/D),\\
 II_u&=D^{-1/2}\sum_{a,s\ {\rm squarefree}\atop Na,Ns>U,\,(a,s)=1}
  \mu(a)g_U(s)\nu(as)\chi_a(u)\chi_s(u)W(NaNs/D).
\end{align*}
Then $M_u=-I_u+II_u$. All ideals here remain outside the same fixed
excluded set. The short $2a_U$ term vanishes because $U<D$ and $W$ is
supported in $[1,2]$. No phase or coprimality mask was removed.

\begin{proposition}[An all-cutoff separate-piece bound is false]
\label{prop:mobius-unit-cutoff}
Take the allowed principal $\nu$, $U=1$, and a fixed nonnegative smooth
$W$ supported in $[1,2]$, bounded below by a positive constant on a closed
subinterval of $(1,2)$ of positive length.
For each of the two pieces,
\[
 \sum_{0<Nu\le H}|I_u|^2\gg_W\frac{D}{\log^2D},\qquad
 \sum_{0<Nu\le H}|II_u|^2\gg_W\frac{D}{\log^2D}
 \quad(H\ge1).
\]
Consequently a proposed bound $H D^\varepsilon$ for each piece at
$H=D^{3/4}$ fails for $\varepsilon<1/4$. In the exact original notation
$H(HD)^\varepsilon$, it fails for $\varepsilon<1/7$.
This refutes an all-cutoff termwise strengthening, not the original signed
mean square or a specified moving-cutoff argument.
\end{proposition}
\begin{proof}
At $U=1$, $a_U=\delta$, so $a_U*a_U*\zeta_0=\zeta_0$ and
$b_U*b_U*\zeta_0=\mu+\zeta_0-2\delta$.
On the annulus the two coefficients are $1$ and $1+\mu(n)$, respectively.
On squarefree ideals the latter is zero or two. At the actual unit row
$u=1$, every original character equals one. A fixed-field prime-ideal
theorem gives $\gg D/\log D$ good prime columns in the subinterval,
so $I_1\gg_W\sqrt D/\log D$. Fix a good prime ideal $q_0$.
Products $q_0p$ with varying distinct good primes in the rescaled
subinterval have $\mu(q_0p)=1$. Their nonnegative contributions give
$II_1\gg_W\sqrt D/\log D$. Keeping that one row in each nonnegative
energy proves the result. The prime-ideal theorem is a classical
unconditional external input, with the field and $q_0$ fixed.
\end{proof}

The factorization does not shorten the primitive row conductor:
for squarefree coprime $a,s$ it is $as$, not $a$ or $s$ separately.
For two resulting columns $n,n'$ with common factor $h$, the oscillatory
conductor is $nn'/h^2$ and a zero-preserving period is $nn'/h$.
The signed low-gcd estimate and the independent positive arithmetic
determinant comparison therefore remain open.

\subsection{Exact factorization, retained zeros and the failed Cauchy route}
\label{sec:mobius-split-details}
\subsubsection{Proof of the regrouping with original coefficients}
The monoid consists of nonzero integral ideals prime to the fixed excluded
set $S$. Every divisor of a good ideal is good, so the identity
$\mu*\zeta_0=\delta$ holds on this same monoid. Expanding convolution
in~\eqref{eq:mobius-split} gives finite divisor sums. Multiply the identity
by $\mu(n)^2\nu(n)\chi_n(u)W(Nn/D)$ and sum. The extra squarefree
indicator does not change $\mu(n)$, since $\mu(n)^3=\mu(n)$.
On its support $n=abc$ forces each factor to be squarefree and any two
to be coprime. Sextic multiplicativity, including zeros, gives
$\chi_n(u)=\chi_a(u)\chi_b(u)\chi_c(u)$; a collision with any factor
makes both sides zero. The fixed character also satisfies
$\nu(n)=\nu(a)\nu(b)\nu(c)$ on these good ideals.

In the small-factor piece group $d=ab$. The remaining variable $c$ is
squarefree and coprime to $d$; $Nd\le U^2$ is necessary, but not sufficient
for a nonzero $\beta_U(d)$: its specified truncated factorizations must
still exist. For each such factorization on squarefree $d$,
$\mu(a)\mu(b)=\mu(d)$, proving the coefficient formula.
In the large-factor piece group $s=bc$. The condition $Nb>U$ implies
$Ns>U$, and summing that divisor gives exactly $g_U(s)$.
Since $Ns>U\ge1$, $s$ is not the unit ideal; its complete M\"obius
divisor sum vanishes, proving the alternative formula for $g_U$.
This regrouping never replaces the two-variable weight by $W(Ns/D)$,
never deletes $(a,s)=1$, and never replaces the resulting coefficients
by a favorable arbitrary family.

One can instead keep the unrestricted convolution identity on nonsquarefree
columns, but then both individual pieces and their local conductor
description change. For $U=1$ and a good prime $p$, the column $p^2$
contributes one to each of $a_U*a_U*\zeta_0$ and
$b_U*b_U*\zeta_0$, even though $\mu(p^2)=0$; these cancel only after
the pieces are recombined. Its squared sextic symbol has a cubic local
component. Thus dropping total squarefreeness is not a harmless
piecewise simplification, although the full unrestricted algebraic
identity still holds.

\subsubsection{The two conductor scales are unchanged}
For a squarefree $n=as$ with $(a,s)=1$, every prime of $n$ has one
nontrivial sextic local component in the row variable. Hence its primitive
finite conductor is the ideal $as$; $NaNs\asymp D$ on the column support.
On multiplying a column by the conjugate of another, the local exponent
is zero at common primes, $1$ at a prime of $n/h$, and $-1$ at a prime
of $n'/h$. The common prime still imposes its original zero on rows.
Writing $n=h r$, $n'=h t$ therefore gives exactly
\[
 \chi_n(u)\overline{\chi_{n'}(u)}
 =\mathbf1_{(u,h)=1}\chi_r(u)\overline{\chi_t(u)},\qquad
 q=rt,\quad\mathfrak l=hrt.
\]
Here $q$ is the primitive oscillatory conductor and $\mathfrak l$ a
period of the full zero extension. These are ideals, not necessarily
squarefree rational norms. In particular, conjugate primes over one split
rational prime are distinct coprime ideal factors. Full defining moduli
of $\nu$ and primary conventions remain fixed in $S$; this row-conductor
calculation does not assert a new uniform reciprocity theorem in the
column variable. In the coprime bulk $Nq\asymp D^2$ remains large.

\subsubsection{Why a first generic Cauchy step does not supply the saving}
On one type-II dyad let $Na\asymp A$, $Ns\asymp B$, $AB\asymp D$.
Keep the full profile and coprimality restrictions in
\[
 X_a(u)=\nu(a)\chi_a(u)
  \sum_{s\ {\rm squarefree}\atop Ns\asymp B,\,(a,s)=1}
  g_U(s)\nu(s)\chi_s(u)W(NaNs/D).
\]
The exact elementary Cauchy inequality is
\[
 \sum_{0<Nu\le H}\left|D^{-1/2}\sum_{Na\asymp A}\mu(a)X_a(u)\right|^2
 \le \frac{\sum_{Na\asymp A}|\mu(a)|^2}{D}
       \sum_{Na\asymp A}\sum_{0<Nu\le H}|X_a(u)|^2.
\]
The $a$-mask has been kept inside $X_a$. Even if one granted an optimistic
uniform slice bound $(H+B)\sum_{Ns\asymp B}|g_U(s)|^2D^\varepsilon$,
the ordinary divisor estimate $|g_U(s)|\le\tau(s)$ and the fixed-field
ideal count would only yield the envelope
\[
 A(H+B)D^\varepsilon\asymp(AH+D)D^\varepsilon.
\]
That hypothetical slice estimate is not established here. This is an
insufficiency of the generic upper-bound calculation, not a lower bound
on the signed energy. At balanced $A,B\asymp D^{1/2}$ and $H=D^{3/4}$,
the $AH$ term loses $D^{1/2}$ relative to the target; the $D$ term also
exceeds $H$. A sharper use of the actual $g_U$, cross-factor constraints
or covariance between the original pieces is required. This calculation
does not refute all bilinear, moving-cutoff or dispersion methods.

The fixed-$U=1$ lower bounds in
Proposition~\ref{prop:mobius-unit-cutoff} are different: they are genuine
energy lower bounds for the separate pieces. If $W\ge w_0>0$ on
$[c_1,c_2]\subset(1,2)$, count primes in
$[c_1D,c_2D]$ for $I_1$ and in
$[c_1D/Nq_0,c_2D/Nq_0]$ for $II_1$. Removing $S$ and $q_0$ removes
finitely many primes. The prime-ideal theorem gives positive constants
times $D/\log D$ in both cases; the latter contribution has coefficient
two. Because every coefficient of $II_1$ is nonnegative, keeping these
products is valid. This is unlike isolating the prime part of the full
signed M\"obius energy, where uncontrolled cross terms remain.
For the exact requested bound $H(HD)^\varepsilon$ at $H=D^{3/4}$,
the competing power is $D^{3/4+7\varepsilon/4}$, giving the threshold
$\varepsilon<1/7$. Replacing $(HD)^\varepsilon$ by $D^\varepsilon$
reparametrizes the loss and gives the equivalent $1/4$ formulation.
The positivity argument is specific to $U=1$. For a concrete control,
the primary element $7-21\omega$ has norm $637=7\cdot7\cdot13$ and
three distinct ideal prime factors. At $U=7$ its type-I coefficient is
$-1$, its type-II coefficient is $-2$, and its M\"obius value is $-1$.
Thus the separate pieces need not have nonnegative coefficients at another
cutoff; no moving-cutoff obstruction follows from the preceding positivity
proof.

\subsubsection{Exact finite checks and formal boundaries}
The deterministic checker represents every good ideal of norm at most
$1024$ once by its primary generator, including inert prime ideals and
squarefree products of three distinct prime ideals.
It checks $1535$ coefficient/regrouping cases at five cutoffs, retains
$19680$ encountered original zeros in actual sextic sums, and verifies
$900$ pair-conductor exponent ledgers. It evaluates principal $\nu$ and
the fixed Kummer character $\nu(n)=\chi_n(2)$ on the same good ideals.
The finite diagnostic weight is the rational profile
$((Nn-D)(2D-Nn))^2/D^4$ on $[D,2D]$, zero elsewhere; it is not claimed
to be the smooth analytic weight of the theorem. Every identity is
algebraic for arbitrary finite weights; these samples do not prove
seminorm or scale-supremum estimates.

For example, at $D=128$, $H=16$, principal $\nu$ and $U=1$, the unit-row
normalized sums have the following exact common scaling:
\[
\begin{array}{c|r}
 D^{9/2}M_1&-122721831\\
 D^{9/2}I_1&314726661\\
 D^{9/2}II_1&192004830
\end{array}
\]
Their exact difference is $II_1-I_1=M_1$. The lower-bound theorem uses
the prime-ideal proof, not extrapolation from these small numbers.
An independent no-production-import oracle also checks convolution,
squarefree regrouping and weighted sextic sums.

\path{MobiusSplit.lean} proves eight actual statements: the inverse-ring
identity, natural-number arithmetic truncation/tail, the M\"obius Dirichlet
identity, the unit cutoff, its exact tail convolution and coefficient,
and nonnegativity. The natural-number Dirichlet ring is an actual
arithmetic object; it is not the ideal convolution ring. Ideal arithmetic,
sextic/Hecke character construction, analytic prime counting and the
unproved signed mean square remain prose or named classical inputs.
No new zero-free region, full GRH proof or positive determinant bridge
has been obtained.


\section{Formal coverage and the next research step}
The new Lean interface checks actual Pl\"ucker polynomials, nodal matrix relations and determinants, finite clock/shift identities and Kraus positivity/trace preservation. The short-row module checks finite-character masks, original finite Gram expansion, matched-row lower bounds and actual Eisenstein norm/evaluation identities. Full planar Poisson summation, ideal-character interfaces, the AR classifications, the derived functors and the complete physics program are not all formalized by these components. Their theorem types and axiom dependencies are printed in the verification logs.

The next load-bearing research task is an estimate for the remaining low-gcd original arithmetic kernel, or an independently constructed positive arithmetic determinant family escaping the precisely stated finite-kernel obstructions. A learned ToricGT-style component may propose a family, choose a search ordering or discover a candidate identity. Acceptance still requires exact algebra, analytic estimates, Lean interfaces where available and independent critical review. A fitted zero picture or a small training loss is an observation, not the desired proof.

% Generated shared reviewed material.
\section{A full-mask Hecke expansion for the type-I piece}
\label{hti:core}
We now continue the signed decomposition through its actual squarefree
inner sum. Fix $K=\mathbb Q(\sqrt{-3})$, the finite-order Hecke character
$\nu$, and the same finite excluded prime set $S$, containing primes above
$6$ and the conductor of $\nu$. For $u\in\mathcal O_K\setminus\{0\}$,
let $\psi_u$ be the primitive Hecke character induced by
$\nu(n)(u/n)_6$ away from $S\cup\operatorname{supp}(u)$.
It is defined before, and independently of, the type-I divisor $d$.
The abelian Kummer extension $K(u^{1/6})/K$ constructs the sextic
factor; multiplying by $\nu$ gives the stated character.
This uses classical abelian reciprocity, not general Artin holomorphy.
Write $\psi_{u,2}$ for the primitive character inducing its square and
$\mathfrak f_u,\mathfrak f_{u,2}$ for the two primitive conductors.
For a squarefree good ideal $d$ put
\[
 T=T(u,d)=S\cup\operatorname{supp}(u)\cup\operatorname{supp}(d),\qquad
 R_T=\prod_{p\in T}Np,\quad Q_u=|D_K|N\mathfrak f_u.
\]
All primes in $T$ are deleted even if absent from either primitive conductor.

\begin{theorem}[Restoration and a uniform scalar contour bound]
\label{hti:theorem}
For $\Re s>1$ the original type-I inner Dirichlet series is
\begin{align}
 F_{u,d}(s)
 &=\sum_{c\ {\rm squarefree}\atop(c,T)=1}
       \frac{\psi_u(c)}{(Nc)^s}
   =\frac{L(s,\psi_u)}{L(2s,\psi_{u,2})}\,C_T(s),\label{hti:ratio}\\
 C_T(s)
 &=\frac{\displaystyle\prod_{p\in T,\ p\nmid\mathfrak f_u}
             (1-\psi_u(p)(Np)^{-s})}
          {\displaystyle\prod_{p\in T,\ p\nmid\mathfrak f_{u,2}}
             (1-\psi_{u,2}(p)(Np)^{-2s})}.
 \label{hti:restore}
\end{align}
It continues meromorphically to $\Re s>1/2$, with only a possible
simple pole at $s=1$ when $\psi_u=1$. In that case its residue is
\[
 \rho_T=\frac{\operatorname*{res}_{s=1}\zeta_K(s)}{\zeta_K(2)}
             \prod_{p\in T}(1+(Np)^{-1})^{-1}.
\]
Fix $0<\delta<1/2$, $\varepsilon>0$ and
$W\in C_c^\infty((0,\infty))$ supported in $[1,2]$. Set
\[\mathcal M_4(W)=\sum_{j=0}^4\sup_x|(x\partial_x)^jW(x)|.\]
Uniformly for $Y\ge1$, $u$ and $d$,
\begin{align}
 \sum_{c\ {\rm squarefree}\atop(c,T)=1}\psi_u(c)W(Nc/Y)
 ={}&\mathbf1_{\psi_u=1}\rho_TY\widehat W(1)\notag\\
 &+O_{\nu,S,\delta,\varepsilon}\!\left(
 \mathcal M_4(W)Y^{1/2+\delta}
 Q_u^{1/4-\delta/2+\varepsilon}R_T^\varepsilon\right).
 \label{hti:scalar}
\end{align}
Here $\widehat W(s)=\int_0^\infty W(x)x^{s-1}\,dx$.
The continuation and functional equation of degree-one Hecke functions
are classical external inputs; the contour argument is given below.
\end{theorem}

At a good prime with $v_p(u)=3$, $\psi_u$ is ramified while its square
is unramified; the second product in~\eqref{hti:restore} must restore a
factor invisible to squaring the zero-extended primitive value.
At $v_p(u)=6$ both characters can forget the prime, but the original row
still vanishes on columns divisible by it. In particular, for principal
$\nu$, $u=\pi^3$ and $u=\pi^6$ at $N\pi=7$, the deleted local factors are
\[
\begin{array}{c|ccc}
 &\text{primitive ratio at }s=1&\text{restoration}&\text{full masked factor}\\\hline
v_p(u)=3&48/49&49/48&1\\
v_p(u)=6&8/7&7/8&1
\end{array}
\]
These are exact local calculations, not numerical evidence about zeros.

For $1\le U$, $U^2\le D$, insert $Y=D/Nd$ into~\eqref{hti:scalar}
and retain the actual $\beta_U(d)$ and outer $\nu(d)\chi_d(u)$.
Classical tame Kummer conductors give $Q_u\ll_\nu Nu$. Thus the normalized
$I_u$ equals its displayed principal-residue sum plus an error $E_u$ with
\begin{equation}\label{hti:outer}
 |E_u|\ll_{\nu,S,\delta,\varepsilon}\mathcal M_4(W)
 D^\delta U^{1-2\delta+\varepsilon}
 (Nu)^{1/4-\delta/2+\varepsilon}.
\end{equation}
The resulting elementary energy envelope is only
\begin{equation}\label{hti:energy}
 \sum_{0<Nu\le H}|E_u|^2
 \ll\mathcal M_4(W)^2D^{2\delta}U^{2-4\delta+\varepsilon}
 H^{3/2-\delta+\varepsilon}.
\end{equation}
It does not prove the short-row target of order $H$ up to small losses.
This is an insufficient upper bound, not a lower bound on the actual
error or a counterexample to coefficient-sensitive averaging. The
original signed mean square and positive determinant bridge stay open.

% Generated shared reviewed material.
\section{Primitive conductors, deleted factors and the scalar contour}
\label{hti:details}
\subsection{Constructing the column character without losing its zeros}
Let $\alpha^6=u$. Since $K$ contains $\mu_6$, the extension $K(\alpha)/K$
is abelian, with character $g\mapsto g(\alpha)/\alpha$. At $p\nmid6u$,
arithmetic Frobenius gives the residue $u^{(Np-1)/6}$, exactly the
source's sextic symbol. Abelian reciprocity makes it a finite-order Hecke
character; multiply by $\nu$ and take its primitive inducing character.
The classical reciprocity input is \cite[Theorem~0.8 and Chapter~I,\S1]{MilneCFTTypeI}.
The finite-order archimedean character is trivial, since $\mathbb C^\times$
is connected.

At $p\nmid6$, write $u=\pi^e v$ with $v$ a local unit. A sixth root of
$v$ lies in an unramified extension: its residue root is separable and
lifts by Hensel's lemma. Inertia acts on a sixth root of $\pi$ through
$\mu_6$, so on $\alpha$ the inertia character has order
$6/\gcd(6,e)$. It is trivial exactly when $6\mid e$; its square is
trivial exactly when $3\mid e$. A nontrivial tame one-dimensional
character has conductor exponent one. Since $\nu$ is unramified outside
$S$, it cannot cancel this inertia there. At primes above $6$, the groups
$K_p^\times/(K_p^\times)^6$ are finite; their finitely many characters
have a uniform maximum conductor. Consequently
\[
 \mathfrak f_u\mid \mathfrak c_\nu
       \prod_{p\mid u,\ p\nmid6}p,\qquad
 Q_u\ll_\nu Nu.
\]
The fixed ideal $\mathfrak c_\nu$ includes the wild and fixed $\nu$ factors.
No assertion identifies $\mathfrak f_u$ with the larger full zero modulus.

For $p\notin T$ put $z=\psi_u(p)(Np)^{-s}$. Its squarefree factor
$1+z=(1-z^2)/(1-z)$ proves the ratio using the \emph{same} deleted set
in numerator and denominator. Restoring the primitive Euler factors yields
\eqref{hti:restore}. At a prime where the primitive square becomes
unramified its value cannot be obtained by squaring the sentinel zero
$\psi_u(p)=0$. For principal $\nu$, $u=\pi^3$ has a quadratic Kummer character
whose primitive square is globally principal. Its local value in that
square is one, not zero. For $u=\pi^6$ both primitive characters are
principal, although the full family still deletes $p$. These two cases
prove the local table directly.

\subsection{Why the half-plane above one half is available}
For a character $\eta$ write $L^T(s,\eta)$ for its Euler product with
all primes of $T$ removed. For $\sigma>1/2$ the product
$1/L^T(2s,\psi_{u,2})$ converges absolutely and is bounded by
$\zeta_K(2\sigma)$, uniformly in the character, $T$ and $t=\Im s$.
Every deleted finite factor is nonzero there. Therefore
\[
 F_{u,d}(s)=L(s,\psi_u)
  \prod_{p\in T,p\nmid\mathfrak f_u}(1-\psi_u(p)(Np)^{-s})
  \big/L^T(2s,\psi_{u,2})
\]
has just the stated possible principal pole. This uses absolute convergence
at $2\sigma>1$, not a conjectural zero-free assertion on the critical line.
For any fixed $\sigma>0$ and $\eta>0$,
\[
 \prod_{p\in T}(1+(Np)^{-\sigma})\ll_{\sigma,\eta}R_T^\eta.
\]
Indeed $1+q^{-\sigma}\le q^\eta$ for all sufficiently large $q$;
the remaining finitely many prime ideals give a uniform fixed constant.
This controls every moving deleted prime; the denominator has already
been bounded by the full convergent zeta product.

We use the classical completed Hecke functional equation in the form
\[
 \Lambda(s,\psi_u)=Q_u^{s/2}(2\pi)^{-s}\Gamma(s)L(s,\psi_u),
 \qquad\Lambda(s,\psi_u)=w_u\Lambda(1-s,\overline{\psi_u}),
 \quad |w_u|=1.
\]
This is a named external theorem, not a newly formalized analytic result;
compare \cite[Theorem~4.3, Corollary~4.4]{KudlaTateTypeI}. The cited notes
sketch some of Tate's calculations; their incomplete displayed epsilon
formulas are not used as independently checked constants here.
At $\Re s=1+\eta$, absolute convergence bounds $L$ by $\zeta_K(1+\eta)$.
The functional equation and Stirling's formula give on $\Re s=-\eta$
the bound $Q_u^{1/2+\eta}(1+|t|)^{1+2\eta}$ up to a fixed constant.
Phragm\'en--Lindel\"of on this strip therefore gives, for fixed
$0<\sigma<1$,
\[
 |L(\sigma+it,\psi_u)|
 \ll_{\sigma,\eta}Q_u^{(1-\sigma)/2+\eta}
                      (1+|t|)^{1-\sigma+2\eta}.
\]
For the principal character apply the strip argument to
$(s-1)L(s,1)/(s+1)$, then divide back on the fixed line $\sigma<1$.
The constants need not be uniform as $\sigma\to1$. Polynomial vertical
strip growth and finite order, as supplied by classical Hecke theory,
justify this application. No Lindel\"of or subconvexity bound is inserted.

\subsection{Mellin inversion, residues and profile uniformity}
Initially, absolute convergence and compact smooth support give
\[
 S_{u,d}(Y)=\frac1{2\pi i}\int_{(1+\eta)}
       \widehat W(s)Y^sF_{u,d}(s)\,ds.
\]
Four integrations by parts using $x\partial_x$ give
$|\widehat W(\sigma+it)|\ll_\sigma\mathcal M_4(W)(1+|t|)^{-4}$.
Choose $0<\eta<\min(\varepsilon/8,1/8)$, decreasing it further
when distributing the small losses, and shift to
$\sigma=1/2+\delta$. The convergent reciprocal Euler product, the finite
factor bound and convexity above control the complete integrand.
Their polynomial $t$ growth is integrable against this Mellin decay and
makes the horizontal edges tend to zero. The only crossed pole is $s=1$
when $\psi_u=1$; evaluating its finite factors gives exactly $\rho_T$.
Taking absolute values on the new vertical line proves
\eqref{hti:scalar}, after renaming the small losses. This also proves
uniformity over any family of profiles with bounded $\mathcal M_4$;
it does not interchange a profile supremum with an arithmetic mean square.

In the original $I_u$, the inner profile is exactly $W(Nc/Y)$ with
$Y=D/Nd$; no $d$-dependent derivative enlargement occurs. On the support
of $\beta_U$, $Nd\le U^2\le D$, hence $Y\ge1$. Also
$R_T\le C_S Nu\,Nd$. The principal term is explicitly
\[
 P_u=\mathbf1_{\psi_u=1}\sqrt D\,\widehat W(1)
 \sum_{d\ {\rm squarefree}\atop(d,S)=1}
       \frac{\beta_U(d)\nu(d)\chi_d(u)}{Nd}\rho_{T(u,d)}.
\]
This sum retains all original phases and zeros; no positivity of it is
assumed. For the error use the actual truncated factorizations, not a free
coefficient bound:
\[
 \sum_{d\ {\rm squarefree}}\frac{|\beta_U(d)|}{(Nd)^{1/2+\delta-\eta}}
 \le\left(\sum_{Na\le U}(Na)^{-1/2-\delta+\eta}\right)^2
 \ll_{\delta,\eta}U^{1-2\delta+2\eta}.
\]
The fixed-field ideal count is $O_K(X)$; dyadic summation proves the last
inequality for sufficiently small $\eta<\delta+1/2$. Inserting this and
the conductor/radical bounds proves~\eqref{hti:outer}.
The elementary lattice count $\#\{u:0<Nu\le H\}=O_K(H)$ then gives
\eqref{hti:energy}. Rows are not divided by units.

At $H=D^{3/4}$ and $U=D^a$ with $0\le a\le1/2$, the power of the
resulting error envelope \emph{relative to} $H$ is
\[
 \frac38+\frac54\delta+(2-4\delta)a
\]
up to the arbitrarily small losses. It remains positive for every
$0<\delta<1/2$. Thus this scalar convexity argument alone supplies no
short-row saving. This is a limitation of the estimate just derived,
not an obstruction to averaging the numerators over Kummer rows, retaining
signed coefficients, or exploiting covariance with type II.

The seminorm uniformity is not a proof of the original scale supremum.
For each independently chosen scale $D'$ the displayed outer argument
requires $U^2\le D'$. If a fixed cutoff is retained while $D'$ shrinks,
then eventually even $D'\le U$: the short $2a_U$ term in the Vaughan
identity can meet the support and must be restored. Neither that term nor
the covariance with type II may be omitted in a supremum or reassembly.

\subsection{Verification boundaries}
\path{HeckeTypeI.lean} checks seven actual complex local/finite-product
statements, including independent primitive coefficient arrays and their
full deleted-factor restoration. It does not identify an ideal Hecke
series with an abstract finite product or formalize a Mellin contour.
The exact checker includes split and inert prime ideals, actual sextic
phases, conductor drops, sixth-power rows and original mask zeros.
The analytic theorem above uses classical abelian reciprocity, Hecke
continuation/functional equation and standard complex-analytic tools as
external inputs. Full ideal-character/formal analytic objects, the original
signed low-gcd mean square and every GRH completion obligation remain open.

